\documentclass[12pt]{article}
\usepackage[margin=1in]{geometry}
\usepackage{titlesec}
\usepackage{hyperref}
\usepackage{enumitem}
\usepackage{listings}
\usepackage{tcolorbox}
\usepackage{xcolor}

% Configure code listings
\definecolor{codebg}{rgb}{0.95,0.95,0.95}
\definecolor{commentcol}{rgb}{0.0,0.5,0.0}
\definecolor{keywordcol}{rgb}{0.0,0.0,0.8}
\definecolor{stringcol}{rgb}{0.8,0.0,0.0}

\lstset{
  backgroundcolor=\color{codebg},
  commentstyle=\color{commentcol},
  keywordstyle=\color{keywordcol},
  stringstyle=\color{stringcol},
  basicstyle=\ttfamily\footnotesize,
  breaklines=true,
  frame=single,
  rulecolor=\color{gray!50},
  xleftmargin=3pt,
  xrightmargin=3pt,
  captionpos=b,
  showstringspaces=false,
  numbers=left,
  numberstyle=\tiny,
  stepnumber=1,
  breakatwhitespace=true,
  tabsize=3,
}

% Define tcolorbox for enhanced code blocks
\tcbuselibrary{breakable,skins}

\newtcolorbox[numbered]{codebox}[1][]{%
  enhanced,%
  colback=codebg,%
  colframe=gray!70,%
  title=#2,%
  #1,%
  left=5pt,right=5pt,top=4pt,bottom=4pt,
  breakable,
}

% Theorem-style environments for code
\usepackage{mdframed}
\mdfsetup{
  skipabove=8pt,
  skipbelow=8pt,
  leftline=true,
  rightline=true,
  topline=true,
  bottomline=true,
  backgroundcolor=codebg,
  linecolor=gray!50,
  innerleftmargin=5pt,
  innerrightmargin=5pt,
  innertopmargin=5pt,
  innerbottommargin=5pt
}

\titleformat{\section}{\large\bfseries\uppercase}{}{0em}{}
\titleformat{\subsection}{\bfseries}{}{0em}{}
\title{MiniANN Project Documentation}
\author{}
\date{}

\begin{document}

\maketitle

\section{Project Overview}
MiniANN is a tiny C++ neural network library for training simple networks (XOR, AND, OR) with a Qt Widgets GUI. The project has three parts: Core Math, Experiment (training pipeline), and GUI.

\section{How to Run the Project}
\paragraph{Method 1: build.bat (recommended)}
\begin{verbatim}
\\ .\\build.bat
\end{verbatim}
Compiles 11 targets (xor_demo.exe through gui_qt.exe). After building:
\begin{verbatim}
\\ .\\gui_qt.exe
\end{verbatim}

\paragraph{Method 2: CMake + Ninja}
\begin{verbatim}
\\ cmake -B build -G Ninja
\\ cmake --build build --target gui_qt
\\ .\\build\\gui_qt.exe
\end{verbatim}

\section{Code Overview — File by File}
This section lists every file you modified or added during this session, with a brief description and the key OOP concepts involved.

\subsection{demos/gui\_qt.cpp — The Qt GUI Window}
\begin{codebox}{Code Block}
\begin{lstlisting}
// Left panel: Dataset buttons (AND, OR, XOR, Iris, CSV),
// config fields (hidden layers, activation, optimizer, epochs, learning rate),
// preset buttons (XOR, Iris, Binary)

// Right panel: Tabs for Train, Validate, Test, Network

// Bottom bar: SAVE MODEL, LOAD MODEL, START TRAINING
// + Speed Slider (1-20, epochs per update tick)
// + [Step 1] button (advances one training step)

// NetWidget class draws the network canvas:
// - Figures out layer sizes from trained network
// - Draws neurons as circles, brightness = activation
// - Draws synapses (edges) between layers
// - Shows source tag at top: "TRAINED — 2 -> 8(tanh)"
// - Hover tooltips: bias, weights, output per neuron

// Key changes in this session:
// 1. Canvas text overlap fix: top inset 22→36px, tag y-position 2→6
// 2. Speed slider + [Step 1] button for incremental training
// 3. QPushButton card styling for START TRAINING
// 4. Left bar formatting: Dataset info + Arch summary in QFrame panels
\end{lstlisting}
\end{codebox}

\textbf{Key functions/sections:}
\begin{itemize}[leftmargin=*]
  \item \textbf{paintEvent()} (line ~562) — Draws canvas background, grid, neurons, edges, source tag, column labels
  \item \textbf{mouseMoveEvent()} (line ~689) — Checks if mouse over neuron → QToolTip with bias/weights/output
  \item \textbf{setNet()} (line ~546) — Called when training finishes; stores network, probe input, tag string
  \item \textbf{describeNet()} (line ~2239) → Converts NeuralNetwork to string like ``2 -> 8(tanh) -> 3(sigmoid)''
  \item \textbf{onTrainClicked()} (line ~2174) — Launches training thread, handles START/STOP
  \item \textbf{m\_speedSlider} + m\_stepBtn — New speed control and step button (lines ~1607-1658)
\end{itemize}

\subsection{include/miniann/optimizer.hpp — Optimizer Strategy Pattern}
\begin{codebox}{Code Block}
\begin{lstlisting}
// Base interface — all optimizers must implement these:
class IOptimizer {
public:
    virtual ~IOptimizer() = default;
    virtual void step(NeuralNetwork& net, std::size_t batchSize) = 0;
    virtual std::string name() const = 0; // runtime identity
};

// Concrete optimizers inherit IOptimizer:
class SGD : public IOptimizer { /* step(), lr() getter */ };
class Adam : public IOptimizer { /* step(), lr(), beta1(), beta2(), eps(), State */ };
class Momentum : public IOptimizer { /* step(), lr(), mu(), State */ };

// Factory — creates optimizer by name without GUI knowing concrete type:
class OptimizerFactory {
public:
    static std::unique_ptr<IOptimizer> create(const std::string& name, double lr);
    static std::unique_ptr<IOptimizer> create(const std::string& name, const OptimizerConfig& cfg);
};

// What changed in this session:
// - Getter methods (lr(), beta1(), beta2(), eps()) added for all three optimizers
// - State structs with state() / restore() for all three (momentum tables, step count)
// - name() overrides returning "sgd", "adam", "momentum" for runtime type ID
\end{lstlisting}
\end{codebox}

\textbf{Key OOP concept:} \textbf{Abstraction + Polymorphism}. The GUI holds \texttt{std::unique\_ptr<IOptimizer>} — it points to a concrete optimizer but only knows the interface type. When you call \texttt{optimizer->step(net, batchSize)}, C++ calls the correct implementation based on the actual object type. Same pattern with \texttt{name()} — status bar always shows ``sgd'', ``adam'', or ``momentum''.

\subsection{include/miniann/loss.hpp — Loss Function System}
\begin{codebox}{Code Block}
\begin{lstlisting}
// Base interface — all loss functions must implement:
class ILoss {
public:
    virtual ~ILoss() = default;
    virtual double compute(const Vector& predicted, const Vector& target) const = 0;
    virtual Vector gradient(const Vector& predicted, const Vector& target) const = 0;
    virtual std::string name() const = 0; // runtime identity
};

// Concrete loss classes:
class MSELoss : public ILoss { /* compute(), gradient(), name() = "mse" */ };
class BCELoss : public ILoss { /* compute(), gradient(), name() = "bce" */ };
class CCELoss : public ILoss { /* compute(), gradient(), name() = "cce" */ };

// What changed in this session:
// - virtual std::string name() const = 0 added to ILoss interface
// - name() overrides in all three classes: "mse", "bce", "cce"
\end{lstlisting}
\end{codebox}

\textbf{Key OOP concept:} \textbf{Abstraction}. The code programs to the \texttt{ILoss} interface, not concrete loss classes. The training pipeline calls \texttt{loss->compute()} and \texttt{loss->gradient()} without knowing MSE/BCE/CCE specifics. The ``name()'' function displays which loss is active in the status bar.

\subsection{src/experiment.cpp — Training Pipeline}
\begin{codebox}{Code Block}
\begin{lstlisting}
// ExperimentController::run(cfg, &g_bridge, &stopRequested) — the main training function
// Runs in a std::thread so GUI doesn't freeze:

g_trainThread = std::make_unique<std::thread>([cfg]() {
    ExperimentResult res = ExperimentController::run(cfg, &g_bridge, &g_bridge.stopRequested);
    // Fills g_bridge with final results under std::lock_guard<std::mutex>
    // g_bridge.hasResult = true, lastNet, lastTrain/Test, full history arrays
    g_bridge.isTraining = false; // training done
});

// Qt timer (~100ms) fires refreshLiveUi() which:
//   Reads g_bridge UNDER std::lock_guard<std::mutex>
//   Updates labels: epoch, train loss, train acc, val acc, test acc
//   Updates plot only if data changed (m_cPlotA/B/V/E gate prevents flicker)

// What fixed in this session:
// - Stratified train/val/test splits: every seed now yields balanced 8/1/1 per class (Iris)
// - Per-feature max-abs normalization (was one global divisor squashing small features to ~0.027)
// - He init for ReLU family (was Xavier, causing seed-dependent collapse to 33%)
// - Preview-call dedup: second refreshPreview() is provably a no-op (fingerprint cache)
\end{lstlisting}
\end{codebox}

\textbf{Key OOP concept:} \textbf{Encapsulation + Threading}. The training thread writes to \texttt{g\_bridge} ``under'' \texttt{std::lock\_guard<std::mutex>}, and the GUI reads ``under'' the same mutex. This prevents two threads from accessing the same data simultaneously (would crash/corrupt). The ``g\_bridge'' struct encapsulates all live/final state, and the mutex is the single access point.

\subsection{build.bat — Build Script}
\begin{codebox}{Code Block}
\begin{lstlisting}
// Step numbering changed from [1/10]...[10/10] to [1/11]...[11/11]
// Added [10/11] Building test_correctness.exe... (mirrors CMake test list)
// Added [11/11] Building gui_qt.exe (Qt Widgets, optional)...

// UCRT g++ auto-selection (prevents MSVCRT fstream crash class)
// deploy_qt.ps1 — run once after building so gui_qt.exe works by double-click
\end{lstlisting}
\end{codebox}

\textbf{Key change:} Step renumbering + test\_correctness.exe addition + URT toolchain auto-detection.

\section{OOP Concepts in This Project — Detailed}
Each concept has a \textbf{Heads-up} summary first, then the technical details with file references.

\subsection{Abstraction (Interfaces)}
\paragraph{Heads-up:} Abstraction means defining a common interface (function signatures) that other classes must follow. Code programs to the interface, not the implementation.

Where: \texttt{include/miniann/optimizer.hpp} :10-15 — \texttt{IOptimizer} with pure virtual \texttt{step()} and \texttt{name()}. Similarly \texttt{include/miniann/loss.hpp} :10-15 — \texttt{ILoss} with virtual \texttt{compute()}, \texttt{gradient()}, \texttt{name()}.

How used: \texttt{SGD}, \texttt{Adam}, \texttt{Momentum} inherit \texttt{IOptimizer} and implement the two pure virtual functions. The GUI holds \texttt{std::unique\_ptr<IOptimizer>} — it points to a concrete optimizer but only knows the interface type. When you call \texttt{optimizer->step(net, batchSize)}, the correct implementation fires based on actual object type (polymorphism). Same for loss functions — GUI doesn't need to know whether MSE, BCE, or CCE is active; it just calls interface methods.

\subsection{Polymorphism (Virtual Functions)}
\paragraph{Heads-up:} Polymorphism means same function name, different behavior depending on actual object type. C++ decides at runtime which version to call.

Where: \texttt{demos/gui\_qt.cpp} — \texttt{g\_bridge.optimizer->step(net, batchSize)} calls correct optimizer's \texttt{step()}. Similarly \texttt{optimizer->name()} returns ``sgd'', ``adam'', or ``momentum'' depending on which optimizer is active.

How used: \texttt{g\_bridge.optimizer} is a \texttt{std::unique\_ptr<IOptimizer>} — could point to \texttt{SGD}, \texttt{Adam}, or \texttt{Momentum}. When \texttt{step()} called, C++ looks at actual object type and runs that optimizer's \texttt{step()}. Analogy: same \texttt{play()} button on different devices — CD player spins disc, Spotify streams, radio turns on. Same button, different behavior. Here, call \texttt{optimizer->step()}, but behavior depends on which optimizer is running.

\subsection{Encapsulation (Data Hiding)}
\paragraph{Heads-up:} Encapsulation means keeping data private inside a class; expose only what's needed through methods. Prevents outside code from messing with internal state.

Where: \texttt{include/miniann/optimizer.hpp} private members + getters (SGD::lr\_ → lr(), Adam::m\_w\_/m\_b\_/t\_ → State struct with state()/restore()). \texttt{g\_bridge.mtx} protects shared state between GUI and training threads. \texttt{NetWidget::m\_tag}, m\_net, m\_probe private with access through setNet().

How used: GUI never touches SGD::lr\_ directly — it calls optimizer->lr() (the getter). Training thread writes to g\_bridge under std::lock\_guard<std::mutex> — GUI reads g\_bridge values safely. NetWidget::m\_hit (hit-test rectangles for tooltips) is private; mouseMoveEvent() checks them but doesn't modify them. Analogy: car — you turn steering wheel and press gas (exposed interfaces), but can't reach in and mess with engine timing belt (hidden inside). Prevents accidental breaking.

\subsection{Factory Pattern}
\paragraph{Heads-up:} Factory pattern creates objects without specifying the exact class type. Ask for ``an optimizer named adam'' and the factory gives you one.

Where: \texttt{include/miniann/optimizer.hpp} :92-97 — OptimizerFactory::create() returns std::unique\_ptr<IOptimizer> based on name string.

How used: 1. GUI combobox selects optimizer name: ``sgd'', ``adam'', or ``momentum''. 2. currentConfig() builds OptimizerConfig from widget values. 2. OptimizerFactory::create(name, config) returns the correct concrete optimizer. 2. GUI stores it in g\_bridge.optimizer (type IOptimizer). 3. Training calls optimizer->step(net, batchSize) — polymorphic call. Analogy: vending machine — press button for ``Adam'', machine gives you Adam optimizer object. You don't need to know how to construct an Adam — the factory handles that. GUI just says ``give me optimizer name X'' and factory delivers it.

\subsection{Composition (Over Inheritance)}
\paragraph{Heads-up:} Composition builds complex things by composing simpler objects, rather than one giant class inheriting from many base classes.

Where: \texttt{NetWidget} has QPainter, QFontMetrics, hit rectangles, tooltips (\texttt{demos/gui\_qt.cpp}). \texttt{ExperimentController} has g\_bridge, g\_trainThread, currentConfig() (\texttt{demos/gui\_qt.cpp}). \texttt{NeuralNetwork} has std::vector\<Layer\>, each Layer has std::vector\<Neuron\> (\texttt{include/miniann/network.hpp}).

How used: \texttt{NetWidget} doesn't inherit from giant NetworkDrawer. Instead, it has a QPainter, has a font metrics object, has a hit-test rectangle vector. It delegates drawing to QPainter, tooltip logic to QToolTip. \texttt{ExperimentController} composes a thread, shared bridge, and config values. \texttt{NeuralNetwork} has-a vector of Layers, each Layer has-a vector of Neurons. Not ``is-a'' relationships. Analogy: LEGO bricks — each brick has one purpose, snap together to make bigger structure.

\subsection{Threading + Mutex (Concurrency)}
\paragraph{Heads-up:} Threading runs code in separate thread so main program doesn't freeze. Mutex (mutual exclusion) prevents two threads from accessing same data at same time (would crash).

Where: \texttt{src/experiment.cpp} — g\_bridge.mtx, g\_trainThread = std::make\_unique<std::thread}(...). GUI reads g\_bridge under mutex, training thread writes g\_bridge under std::lock\_guard<std::mutex>.

How used: \texttt{onTrainClicked()} launches std::thread running ExperimentController::run(). Training thread writes metrics to g\_bridge — always under std::lock\_guard<std::mutex>. \texttt{refreshLiveUi()} (Qt timer, ~100ms) reads g\_bridge — also under std::lock\_guard<std::mutex>. Ensures: while training thread writes ``epoch 42, loss 0.123'', GUI isn't in middle of reading epoch 42's data → no garbled values, no crashes. Analogy: bathroom with ``Occupied'' sign — when someone's inside (training thread), sign is up (mutex locked). GUI (waiting outside) can't enter until person inside finishes and leaves (releases mutex). Prevents two people in bathroom at same time.

\section{Execution Trace — What Happens When You Run It}
\paragraph{Step 0: Application starts} → gui\_qt.exe launches, QApplication created, main window shown, left panel config visible, right panel tabs (Train/Validate/Test/Network), bottom: SAVE/LOAD/START + Speed Slider + [Step 1].

\paragraph{Step 1: User configures and clicks [START TRAINING]} → User sets XOR, 2 hidden layers, Tanh, Adam, 1500 epochs, lr 0.01 → clicks START → onTrainClicked() fires.

\paragraph{Step 2: Training runs in background} → onTrainClicked() launches std::thread running ExperimentController::run(cfg, \&g\_bridge, \&g\_bridge.stopRequested). Thread iterates epochs: forward pass → compute loss (MSE/BCE/CCE) → compute gradients → optimizer->step(net, batchSize) → update weights → evaluate accuracy → write live data to g\_bridge UNDER mutex → check stopRequested.

\paragraph{Step 3: GUI reads live data} → Every ~100ms Qt timer fires refreshLiveUi(). Reads g\_bridge UNDER mutex → updates labels (epoch, train loss, train acc, val acc, test acc) → updates plot only if data changed (m\_cPlotA/m\_cPlotB/m\_cPlotV/m\_cPlotE gate prevents flicker on idle).

\paragraph{Step 4: Training completes} → ExperimentController::run() exits thread → g\_bridge.hasResult=true, g\_bridge.lastCompleted=true/false, g\_bridge.lastNet=trained pointer, full history populated → std::thread.join() in onTrainClicked() → GUI resumes.

\paragraph{Step 5: GUI updates to show results} → refreshConfigUi() runs → metric labels set → status bar badge (COMPLETED/STOPPED/ERROR) → m\_netview->setNet(net, probe, tag) → tag = ``TRAINED — 2 -> 8(tanh)'' (from describeNet()) → speed slider respected for render density → [Step 1] button can advance one step.

\paragraph{Step 6: User probes neurons} → mouseMoveEvent() checks m\_hit rectangles → if inside → QToolTip::showText() with: input neuron → ``Input x1 = 0.5''; hidden/output neuron → ``L1N2 b=0.35 out=0.82\nw=[0.12, -0.45, 0.91]''.

\paragraph{Step 7: User steps forward incrementally} → [Step 1] button launches single-epoch training thread (same code path as START, but only one run completes, then thread exits).

\paragraph{Step 8: Save/Load model} → SAVE MODEL → serialize NeuralNetwork to .model file (layer sizes, weights, activations) → LOAD MODEL → browse to .model → network appears in GUI, ready for Network tab or further training.

\paragraph{Step 9: Switch tabs} → Train tab (metrics plot), Validate tab (validation metrics), Test tab (test metrics + confusion matrix), Network tab (network diagram with live neuron firing).

\section{Quick Reference — Where Everything Is}
\begin{tabular}{lll}
\textbf{Concern} & \textbf{File} & \textbf{Line/Section} \\
OOP Abstraction & include/miniann/optimizer.hpp & :10-15 (IOptimizer) \\
OOP Polymorphism & demos/gui\_qt.cpp & g\_bridge.optimizer->step() \\
OOP Encapsulation & include/miniann/optimizer.hpp & :19-56 (private + getters) \\
OOP Factory & include/miniann/optimizer.hpp & :92-97 (OptimizerFactory) \\
OOP Loss Interface & include/miniann/loss.hpp & :10-15 (ILoss) \\
Training Thread & src/experiment.cpp & :2188-2236 (g\_trainThread) \\
Shared State (g\_bridge) & src/experiment.cpp + global & Mutex-protected \\
NetWidget Canvas & demos/gui\_qt.cpp & :562-687 (paintEvent) \\
Speed Slider & demos/gui\_qt.cpp & :1607-1658 (m\_speedSlider) \\
Canvas Text Fix & demos/gui\_qt.cpp & :568, :684 (inset 22→36px) \\
Build Script & build.bat & :36-76 (11/11 numbering) \\
Qt Deployment & deploy\_qt.ps1 & Run once after building \\
\end{tabular}

\section{FAQ — Common Questions}
\begin{description}
\item[Q: Where is OOP used in this project?]
A: 1. include/miniann/optimizer.hpp — IOptimizer base class with virtual step() and name(); SGD, Adam, Momentum inherit and implement. 2. include/miniann/loss.hpp — ILoss base class with virtual compute(), gradient(), name(); MSELoss, BCELoss, CCELoss inherit. 3. demos/gui\_qt.cpp — NetWidget composes QPainter, QFontMetrics, hit rectangles, tooltips (composition). 4. src/experiment.cpp — g\_bridge mutex-protected shared state read by GUI, written by thread; g\_trainThread for background training. 5. include/miniann/optimizer.hpp — OptimizerFactory::create() returns std::unique\_ptr<IOptimizer> based on name. 6. demos/gui\_qt.cpp — onTrainClicked() launches std::thread; currentConfig() builds OptimizerConfig; factory creates optimizer.

\item[Q: How does training actually happen when I run gui\_qt.exe?]
A: 1. Launch app → main window appears. 2. Configure settings (dataset, layers, activation, optimizer, epochs) and click START TRAINING. 3. onTrainClicked() launches background std::thread running ExperimentController::run(). 4. Thread trains for specified epochs, writing live metrics (loss, accuracy) to shared g\_bridge struct (protected by mutex). 5. Every ~100ms, Qt timer fires refreshLiveUi(), reads g\_bridge under mutex, updates labels and plot. 6. When training finishes, thread joins back, GUI displays final metrics, shows network diagram, enables Save/Load model.

\item[Q: What does the Speed Slider do?]
A: Controls epochs rendered per GUI update tick. Value 1 = render every epoch (fastest visible). Value 20 = render 1 epoch every 20 actual epochs (slow-mo, helpful for watching learning on XOR/AND/OR). [Step 1] button advances training by one step, respecting speed setting.

\item[Q: Why was the TRAINED tag cutting through neurons?]
A: Tag drawn at area.top() + 2 inside frame with 22px top inset. First neuron row at area.top() + 18, leaving only 16px for tag text plus DPI scaling. Fix: increased top inset from 22px to 36px (6px above tag, 20px between tag and first neuron row), moved tag y-position from area.top() + 2 to area.top() + 6.

\item[Q: Can I add more datasets?]
A: Yes — AND, OR, XOR, Iris, CSV buttons in left panel. Add more by adding buttons, data generation code, possibly adjust boundary visualization for 2D inputs.

\item[Q: How do I save and load a trained network?]
A: SAVE MODEL → serialize NeuralNetwork to .model file (layer sizes, weights, activations). LOAD MODEL → browse to .model → network appears in GUI, ready for Network tab or further training.

\item[Q: What if I want a different optimizer or loss?]
A: Use dropdown comboboxes in left config panel. Optimizer: SGD, Adam, Momentum. The OptimizerFactory creates the right object by name. ILoss interface ensures all losses have same API (compute(), gradient(), name()).

\item[Q: Where can I learn more about the math inside?]
A: include/miniann/ : network.hpp (NeuralNetwork, Layer, Neuron structs, Vector type), optimizer.hpp (how step() updates weights), loss.hpp (how compute() and gradient() calculate error), experiment.hpp (training pipeline and result types), src/*.cpp (actual implementations).
\end{description}

\end{document}